% SPDX-License-Identifier: CC-BY-SA-4.0
% Author: Matthieu Perrin
% Part: 
% Section: 
% Sub-section: 
% Frame: 

\begingroup

\begin{frame}[fragile]{Démonstration de la linéarisabilité}

  \begin{lstlisting}[numbers=left, gobble=4]
    public void push(T value) {
      int i = range.getAndIncrement();
      items.set(i, value);
    }
    public T pop() {
      for(int i = range.get()-1; i>=0; i--) {
        T value = items.getAndSet(i, null);
        if(value != null) return value;
      }
      return null;
    }
  \end{lstlisting}
  
  Soit $range(v_i)$ la valeur de \lstinline{range.getAndIncrement()} lors de $push(v_i)$.

  On définit une relation $\le$ sur les opérations d'une exécution :
  \begin{itemize}
  \item $\structure{push(v)} < \structure{pop()\rightarrow v}$
  \item Si $\structure{range(v_1)} < \example{range(v_2)}$
    \begin{itemize}
    \item si $\example{pop() \rightarrow v_2}$ a lu $\structure{items[range(v_1)]}$ après qu'elle a été écrite
      \begin{itemize}
      \item alors $\structure{push(v_1)} < \structure{pop()\rightarrow v_1} < \example{push(v_2)} < \example{pop()\rightarrow v_2}$
      \end{itemize}
    \item Sinon
      \begin{itemize}
      \item $\structure{push(v_1)} < \example{push(v_2)} < \example{pop()\rightarrow v_2} < \structure{pop()\rightarrow v_1}$
      \end{itemize}
    \end{itemize}
  \item Intercaller les opérations $pop()\rightarrow \bot$ et $push(v)$ restantes
  \end{itemize}
  On démontre que $\le$ est une relation d'ordre totale, qui contient le temps réel. 
  
\end{frame}

\endgroup
\endinput
